\documentclass[12pt,dvipdfmx]{beamer}
\usepackage{pgfpages}
\usepackage{graphicx}
\usepackage{listings,jlisting}
\usepackage{fancybox}
\usepackage{hyperref}
\usepackage{multimedia}

\graphicspath{{out/tex/svg/}}

%%%%%%%%%%%%%%%%%%%%%%%%%%%
%%% themes
%%%%%%%%%%%%%%%%%%%%%%%%%%%
\usetheme{Rochester}
%% no navigation bar
% default boxes Bergen Boadilla Madrid Pittsburgh Rochester
%% tree-like navigation bar
% Antibes JuanLesPins Montpellier
%% toc sidebar
% Berkeley PaloAlto Goettingen Marburg Hannover Berlin Ilmenau Dresden Darmstadt Frankfurt Singapore Szeged
%% Section and Subsection Tables
% Copenhagen Luebeck Malmoe Warsaw

%%%%%%%%%%%%%%%%%%%%%%%%%%%
%%% innerthemes
%%%%%%%%%%%%%%%%%%%%%%%%%%%
% \useinnertheme{circles}	% default circles rectangles rounded inmargin

%%%%%%%%%%%%%%%%%%%%%%%%%%%
%%% outerthemes
%%%%%%%%%%%%%%%%%%%%%%%%%%%
% outertheme
% \useoutertheme{default}	% default infolines miniframes smoothbars sidebar sprit shadow tree smoothtree


%%%%%%%%%%%%%%%%%%%%%%%%%%%
%%% colorthemes
%%%%%%%%%%%%%%%%%%%%%%%%%%%
\usecolortheme{seahorse}
%% special purpose
% default structure sidebartab 
%% complete 
% albatross beetle crane dove fly seagull 
%% inner
% lily orchid rose
%% outer
% whale seahorse dolphin

%%%%%%%%%%%%%%%%%%%%%%%%%%%
%%% fontthemes
%%%%%%%%%%%%%%%%%%%%%%%%%%%
\usefonttheme{serif}  
% default professionalfonts serif structurebold structureitalicserif structuresmallcapsserif

%%%%%%%%%%%%%%%%%%%%%%%%%%%
%%% generally useful beamer settings
%%%%%%%%%%%%%%%%%%%%%%%%%%%
% 
\AtBeginDvi{\special{pdf:tounicode EUC-UCS2}}
% do not show navigation
\setbeamertemplate{navigation symbols}{}
% show page numbers
\setbeamertemplate{footline}[frame number]


%%%%%%%%%%%%%%%%%%%%%%%%%%%
%%% define some colors for convenience
%%%%%%%%%%%%%%%%%%%%%%%%%%%

\newcommand{\mido}[1]{{\color{green}#1}}
\newcommand{\mura}[1]{{\color{purple}#1}}
\newcommand{\ore}[1]{{\color{orange}#1}}
\newcommand{\ao}[1]{{\color{blue}#1}}
\newcommand{\aka}[1]{{\color{red}#1}}

\setbeamercolor{syntax}{bg=cyan!20!white}
\setbeamercolor{example}{bg=yellow!20!white}
\setbeamercolor{output}{bg=white}

%%%%%%%%%%%%%%%%%%%%%%%%%%%
%%% how to typset code
%%%%%%%%%%%%%%%%%%%%%%%%%%%

\lstset{language = python,
numbers = left,
numberstyle = {\tiny \emph},
numbersep = 10pt,
breaklines = true,
breakindent = 40pt,
frame = tlRB,
frameround = ffft,
framesep = 3pt,
rulesep = 1pt,
rulecolor = {\color{blue}},
rulesepcolor = {\color{blue}},
flexiblecolumns = true,
keepspaces = true,
basicstyle = \ttfamily\small,
identifierstyle = ,
commentstyle = ,
stringstyle = ,
showstringspaces = false,
tabsize = 4,
escapechar=\@,
xrightmargin=3zw,
}


\title{微分方程式・偏微分方程式}
\institute{東京大学}
\author{田浦健次朗 \\ 電子情報工学科}
\date{}

\AtBeginSection[] % Do nothing for \section*
{
\begin{frame}
\frametitle{Contents}
\tableofcontents[currentsection,currentsubsection]
\end{frame}
}

\begin{document}
\maketitle

%%%%%%%%%%%%%%%%%%%%%%%%%%%%%%%%%% 
\begin{frame}
\frametitle{Contents}
\tableofcontents
\end{frame}


\iffalse
%%%%%%%%%%%%%%%%% %%%%%%%%%%%%%%%%% 
\begin{frame}
\frametitle{表記上の慣習}

\begin{itemize}
\item \ao{$t$} : 時刻 (実数)
\item \ao{$x$} : 場所 (シミュレートしたい次元に応じて実数または2,3次元のベクトル)
\item \ao{$u$} : 未知の関数
\item \ao{$f, g, \ldots$} : 与えられた関数
\end{itemize}
\end{frame}
\fi


\section{微分方程式}
%%%%%%%%%%%%%%%%% %%%%%%%%%%%%%%%%% 
\begin{frame}
\frametitle{微分方程式とは}
\begin{itemize}
\item 未知の関数(例えば$x$に関する未知の関数$y(x)$)に関する方程式で,
  $x$, $y(x)$, $y'(x)$, $y''(x)$, \ldots などを含んだ式
  \begin{itemize}
  \item しばしば$y(x)$の$(x)$を省略して単に$y$と書く
  \item $y$が1変数のとき, \ao{「常」}微分方程式
  \item 方程式の\ao{「階数」}は$y$の最大の微分の回数($'$の個数)
  \item 例 (1階の常微分方程式)
    \[ y' - x y = 0 \]
  \end{itemize}
\item 微分方程式を「解く」とは?
  \begin{itemize}
  \item $y$を$x$の関数として求めること
  \item 例 (上記の解は実は)
    \[ y = Ce^{\frac{1}{2}x^2} \;\; \mbox{($C$は定数)} \]
  \item どう求めたかは重要でないのでスキップ
  \end{itemize}
\end{itemize}
\end{frame}


%%%%%%%%%%%%%%%%% %%%%%%%%%%%%%%%%% 
\begin{frame}
\frametitle{初期値}
\begin{itemize}
\item 通常, 
  \[ y' - x y = 0 \]
  のような式だけでは$y$は一意に決まらない
\item この解は以下($C$は任意の定数)だった
  \[ y = Ce^{\frac{1}{2}x^2} \;\; \mbox{($C$は定数)} \]
\item そこで通常,
  ある$x=a$での$y$の値$y(a)=c$ \ao{(初期値)}を一緒に与えて解く
\item 例
  \[ y(0) = 1 \]
  という初期値を与えると,
  \[ y = e^{\frac{1}{2}x^2} \]
\end{itemize}
\end{frame}

%%%%%%%%%%%%%%%%% %%%%%%%%%%%%%%%%%
\begin{frame}
\frametitle{記号的な「解く」と数値計算での「解く」}
\begin{itemize}
\item 記号的な求解 $=$
  $y(x)$を$x$の「式」として表す
  \begin{itemize}
  \item 通常, 数学の問題で解くといえばこちら
  \end{itemize}
  
\item 数値的な求解 $=$
  \ao{$y(x)$の近似値を任意の$x$に対して求める}
  \begin{itemize}
  \item 通常, コンピュータで解くといえばこちら
  \end{itemize}
\end{itemize}
\end{frame}

\section{常微分方程式の数値的な求解}
%%%%%%%%%%%%%%%%% %%%%%%%%%%%%%%%%% 
\begin{frame}
\frametitle{常微分方程式の数値的な求解}
\begin{itemize}
\item 以下の形($y' = \ldots$)の1階の常微分方程式
\begin{eqnarray*}
\frac{dy}{dx} & = & f(x, y), \\
y(a) & = & c
\end{eqnarray*}
\begin{itemize}
\item $y$ : 未知の, $x$の関数 ($y(x)$)
\item $a, c$ : 与えられた定数 (初期値)
\item $f$ : 与えられた関数
\end{itemize}

\item<2-> 基本アイデア:
  \begin{itemize}
  \item ある$x$における$y$の値がわかれば右辺が計算できる
  \item $\rightarrow$ その時点での$\frac{dy}{dx}$がわかる
  \item $\rightarrow$ 少し異なる$x$に対する$y$ (つまり$y(x + h)$)がわかる
  \end{itemize}

\item<3-> 上記を形式的に実行する方法 $\rightarrow$ \ao{離散化}
  
\end{itemize}
\end{frame}

%%%%%%%%%%%%%%%%% %%%%%%%%%%%%%%%%% 
\begin{frame}
\frametitle{離散化}
\begin{itemize}
\item 数式で書けば: $h$を十分小さく取れば，
\begin{equation}
\frac{y(x + h) - y(x)}{h} \approx f(x, y(x)) \label{eq:ode}
\end{equation}
\[ \therefore y(x+h) \approx y(x) + f(x, y(x)) h \]

\item []
\[ y(a) \rightarrow y(a+h) \rightarrow y(a+2h) \rightarrow \cdots \]

\item キーは式(\ref{eq:ode})で，元の微分方程式の\ao{「離散化」}という
\end{itemize}
\end{frame}

%%%%%%%%%%%%%%%%% %%%%%%%%%%%%%%%%% 
\begin{frame}[fragile]
\frametitle{プログラム化}
\begin{itemize}
\item []
  \begin{eqnarray*}
    y(a+h) & \approx & y(a) + f(a, y(a)) h \\
    y(a+2h) & \approx & y(a+h) + f(a+h, y(a+h)) h \\
    y(a+3h) & \approx & y(a+2h) + f(a+2h, y(a+2h)) h \\
    & \ldots & 
  \end{eqnarray*}
\item 大雑把には,
\begin{lstlisting}
繰り返す:
    y = y + f(x, y) * h
    x = x + h
\end{lstlisting}
\end{itemize}
\end{frame}

%%%%%%%%%%%%%%%%% %%%%%%%%%%%%%%%%% 
\begin{frame}[fragile]
\frametitle{ちゃんとプログラム化}
\begin{itemize}
\item []

\begin{lstlisting}
# 初期値    
x = a
y = c
h = (b - a) / n_steps   # 刻み幅
for i in range(n_steps):
    y = y + f(x, y) * h
    x = x + h
# この時点で y @$\approx y(b)$@
\end{lstlisting}

\item 誤差の積もりにくい計算の仕方(離散化の方法)とかありますが，
  基本はこれだけ!
\end{itemize}
\end{frame}

%%%%%%%%%%%%%%%%% %%%%%%%%%%%%%%%%% 
\begin{frame}[fragile]
  \frametitle{シミュレーションでよくある状況: $x$が時刻で$y$が位置}
\begin{itemize}
\item 時刻は$t$, 位置は$u$で表すことにする

\item 微分方程式 : 時刻と位置 $\rightarrow$ 速度
  \[ \frac{du}{dt} = f(t, u) \]
\end{itemize}
\begin{columns}
    \begin{column}{0.45\textwidth}
\begin{itemize}
\item 離散化 : ある時刻における位置 $\rightarrow$ 一瞬後の位置
  \[ u(t + \Delta t) \approx u(t) + f(t, u) \Delta t \]
\end{itemize}
\end{column}
\begin{column}{0.1\textwidth}
\end{column}
\hspace{5mm}
    \begin{column}{0.6\textwidth}
\begin{lstlisting}
t = a
u = c
dt = (b - a) / n_steps
for i in range(n_steps):
    # 位置+速度@$\times$@微小時間
    u = u + f(t, u) * dt
    t = t + dt
# u @$\approx$@ 時刻bにおける位置    
\end{lstlisting}
\end{column}

  \end{columns}
\end{frame}

%%%%%%%%%%%%%%%%% %%%%%%%%%%%%%%%%% 
\begin{frame}[fragile]
\frametitle{2階だったら?}
\begin{itemize}
\item 力学の多くの微分方程式($ma = F$)は,
\begin{eqnarray*}
\mbox{加速度} & = & f(\mbox{時刻}, \mbox{位置}, \mbox{速度}) \\
\ddot{x} & = & f(t, x, \dot{x})
\end{eqnarray*}

\item これのプログラムは
\begin{lstlisting}
t = a
u = .. # 初期位置
v = .. # 初速度
dt = (b - a) / n_steps
for i in range(n_steps):
    u = u + v * dt
    v = v + f(t, u, v) * dt
    t = t + dt
\end{lstlisting}
\end{itemize}
\end{frame}

\iffalse
%%%%%%%%%%%%%%%%% %%%%%%%%%%%%%%%%%
\begin{frame}
\frametitle{復習: (常)微分方程式をプログラムで解くとは}
\begin{itemize}
\item [] 
\begin{eqnarray*}
\frac{du}{dt} & = & f(u, t), \\
u(a) & = & c
\end{eqnarray*}
\item 解釈:
  \begin{itemize}
  \item ある時点$t$における$u$がわかれば右辺が計算できる
  \item $\rightarrow$ その時点での$\frac{du}{dt}$がわかる
  \item $\rightarrow$ 一瞬あとの$u$ (つまり$u(t+\Delta t)$)がわかる
  \end{itemize}
\end{itemize}
\end{frame}

%%%%%%%%%%%%%%%%% %%%%%%%%%%%%%%%%% 
\begin{frame}
\frametitle{離散化}
\begin{itemize}
\item 数式で書けば: $\Delta t$を十分小さく取れば，
\begin{equation}
\frac{u(t+\Delta t) - u(t)}{\Delta t} \approx f(u(t), t) \label{eq:ode}
\end{equation}
\[ \therefore u(t+\Delta t) \approx u(t) + f(u(t), t) \Delta t \]

\item []
\[ u(a) \rightarrow u(a+\Delta t) \rightarrow u(a+2\Delta t) \rightarrow \cdots \]

\item キーは式(\ref{eq:ode})で，元の微分方程式の\ao{「離散化」}という
\end{itemize}
\end{frame}


%%%%%%%%%%%%%%%%% %%%%%%%%%%%%%%%%% 
\begin{frame}[fragile]
\frametitle{プログラム化}
\begin{itemize}
\item [] 
\[ u(t+\Delta t) \approx u(t) + f(u(t), t) \Delta t \]
\item プログラムにすると,

\begin{lstlisting}
t = a
u = c
dt = (b - a) / n_steps
for i in range(n_steps):
    u = u + f(u, t) * dt
    t = t + dt
\end{lstlisting}

\item 誤差の積もりにくい計算の仕方(離散化の方法)とかありますが，
  基本はこれだけです!
\end{itemize}
\end{frame}


%%%%%%%%%%%%%%%%% %%%%%%%%%%%%%%%%% 
\begin{frame}[fragile]
\frametitle{一応, 2階の場合も}
\begin{itemize}
\item [] 
  \begin{eqnarray*}
\frac{d^2u}{dt^2} & = & f(u, u', t) \\
u(a) & = & c, \\
u'(a) & = & d
  \end{eqnarray*}

\item []

\begin{lstlisting}
t = a
u = c   # u(t)
u_ = d  # u'(t)
dt = (b - a) / n_steps
for i in range(n_steps):
    u  = u  + u_ * dt
    u_ = u_ + f(u,u_,t) * dt
    t = t + dt
\end{lstlisting}

\end{itemize}
\end{frame}
\fi


%%%%%%%%%%%%%%%%% %%%%%%%%%%%%%%%%% 
\begin{frame}[fragile]
\frametitle{感染拡大の数理モデル}
\begin{itemize}
\item 人口を以下に分類 ($t$ : 時刻)
  \begin{itemize}
  \item $S(t)$(useptible) : これから感染するかもしれない人(まだ陰性)
  \item $I(t)$(nfected) : 感染中
  \item $R(t)$(ecovered) : 治った
  \end{itemize}
\item 方程式
\[\begin{array}{rcll}
    S'(t) & = & -\beta S I & \mbox{$S$と$I$が出会うと一定確率で感染} \\
    I'(t) & = & \beta S I - \gamma I \\
    R'(t) & = & \gamma I & \mbox{感染者は一定時間に一定確率で治る} \\
  \end{array}\]
\item 注: $S' + I' + R' = 0$, すなわち$S + I + R =$一定
\end{itemize}
\end{frame}


%%%%%%%%%%%%%%%%% %%%%%%%%%%%%%%%%% 
\begin{frame}[fragile]
\frametitle{感染拡大の数理モデル}
\begin{itemize}
\item プログラムも簡単
\item ちなみに,
  \[ I'(t) = (\beta S - \gamma) I \]
  なので,
  \[ \beta S - \gamma > 0 \]
  つまり
  \[ \frac{\beta S}{\gamma} > 1 \]
  だと感染拡大 (左辺 $=$ 基本再生産数)
\end{itemize}
\end{frame}

\iffalse
%%%%%%%%%%%%%%%%% %%%%%%%%%%%%%%%%% 
\section{(常)微分方程式を解くもう一つのパターン}
%%%%%%%%%%%%%%%%% %%%%%%%%%%%%%%%%% 
\begin{frame}
  \frametitle{(常)微分方程式を解くもう一つのパターン}

  \begin{itemize}
  \item 微分方程式を解くのに一番簡単なパターンは,
    \begin{itemize}
    \item 初期状態(位置と速度)が与えられ,
    \item その状態から少し後の状態がすべて計算できる
    \end{itemize}
    のでそれを繰り返せば先の状態がわかる, というパターン
    
  \item 別のパターンは,
    \begin{itemize}
    \item 初期値の一部(例えば位置だけ)と,
    \item 目標(たとえばある時刻 $T$ における位置)
    \end{itemize}
    が与えられ, 初期値の残り(例えば初速度)を求めるというパターン
  \end{itemize}
\end{frame}


%%%%%%%%%%%%%%%%%%%%%%%%%%%%%%%%%% 
\begin{frame}
  \frametitle{例題}
  \begin{itemize}
  \item $(0,0)$から玉を発射して, 時刻1で$(1,0)$にいる的に命中させたい
  \item 玉は重力と, 速度に比例する抵抗を受けて進む
    \[ ma = mg - kv \]
  \item そのためには初速をいくらにしたらよいか?
  \item 数学的に言うと, わかっているのは, 内部の点における微分方程式と,
    「端の点」における値($x(0), x(1)$)であるということで,
    境界値問題と言われる
  \end{itemize}
\end{frame}


%%%%%%%%%%%%%%%%%%%%%%%%%%%%%%%%%% 
\begin{frame}
  \frametitle{他の例題}
  \begin{itemize}
  \item このような問題は, 「微分」が時間方向の微分ではなく,
    空間方向の微分であるという場合に, もっとよく現れる
  \item 例:
    \begin{itemize}
    \item 鉄の棒があって, 左端の温度が$T_0$,
      右端の温度が$T_1$に保たれている. さてその途中はどんな温度になるか?
    \item 重さのある柔らかい紐の端点を適当に固定する.
      紐はどんな曲線を描くか (懸垂線)?
    \end{itemize}

  \item 2次元以上では更に良く現れ, 重要
    \begin{itemize}
    \item 鉄の板があって, 縁の温度が与えられている.
      内部はどんな温度になるか?
      (内部では, 「熱伝導の方程式」が満たされている)
    \item 平行平板コンデンサが周りに作る電位はどんな電場か?
      (一方の電極の電位$=0$, もう一方の電極の電位が$=1$という
      条件で, それ以外の場所の電位を求めるという問題.
      電位が満たす方程式は「マックスウェルの方程式」から導かれる)
    \end{itemize}
  \end{itemize}
\end{frame}

%%%%%%%%%%%%%%%%%%%%%%%%%%%%%%%%%% 
\begin{frame}
  \frametitle{基本アイデア}

  軌跡上に多数の点をとり,
  それらの間に(運動方程式に従って)成り立つはずの,
  連立方程式を作り, それを解く
  
  \begin{itemize}
  \item 時間$[0,1]$を$n$等分し, $t_0, t_1, ..., t_n$ とする ($t_i = i/n$)
  \item $t_i$における位置を$(x_i, y_i)$と書き$x_0, y_0, x_n, y_n$はgiven.
    $x_1, \ldots , x_{n-1}, y_1, \ldots , y_{n-1}$が未知数),
    運動方程式を元にこれらの間に成り立つ方程式を立てる
    
  \item つまり, 運動方程式を離散化する
  \end{itemize}
\end{frame}


%%%%%%%%%%%%%%%%%%%%%%%%%%%%%%%%%% 
\begin{frame}
  \frametitle{運動方程式の離散化}

    \[ ma = mg - kv \]
    \[ \ddot{x}(t) = g - \frac{k}{m} \dot{x}(t) \]

    \begin{itemize}
    \item
      \[ \dot{{x}}(t) \approx \frac{{x(t+h) - x(t)}}{{h}} \]
    \item
      \[ \ddot{{x}}(t) \approx \frac{{\dot{{x}}(t+h) - \dot{{x}}(t)}}{{h}} = \frac{{x(t+2h) - 2x(t+h) + x(t)}}{{h^2}}  \]
      
    \item
      \[ \frac{{x_{{i+2}} - 2 x_{{i+1}} + x_i}}{{h^2}} =     - \frac{{k}}{{m}} \frac{{x_{{i+1}} - x_i}}{{h}} \]
      
    \item 
      \[ \frac{{y_{{i+2}} - 2 y_{{i+1}} + y_i}}{{h^2}} = - g - \frac{{k}}{{m}} \frac{{y_{{i+1}} - y_i}}{{h}} \]
    \item \[ (i = 0, \ldots, n - 2) \]
    \end{itemize}
\end{frame}
\fi

%%%%%%%%%%%%%%%%% %%%%%%%%%%%%%%%%% 
\section{偏微分方程式}
%%%%%%%%%%%%%%%%% %%%%%%%%%%%%%%%%% 

%%%%%%%%%%%%%%%%% %%%%%%%%%%%%%%%%% 
\begin{frame}
\frametitle{まだできないこと}
\begin{itemize}
\item これでシミュレートできるのはせいぜい数個の質点
(または数個のパラメータ(e.g., 角度)だけで決まる系)の動き

\item 物理の多くの問題で求めたいものは，「時刻」と「場所」の関数\ao{(「場」)}
  \begin{itemize}
  \item 電場，磁場，重力場，鉄板表面の温度，速度の分布，圧力の分布，力の分布，
    電子の波動関数(存在確率), \ldots
  \item 記号で書けば（$u(t)$ではなく），\ao{$u(x, y, z, t)$}
  \end{itemize}
\item 別の言い方をすると，\ao{時間の関数が無数に($x, y, z$ごとに)ある}
\end{itemize}
\end{frame}



%%%%%%%%%%%%%%%%% %%%%%%%%%%%%%%%%% 
\begin{frame}
\frametitle{場を記述する方程式}
\begin{itemize}
\item 「時間($t$)」と「場所($x, y, z$)」の関数$u(x, y, z, t)$の方程式
\item または定常状態を求める問題では場所だけの関数$u(x, y, z)$のこともある
\item 式としては$\frac{\partial u}{\partial t}$,
$\frac{\partial u}{\partial x}$,
$\frac{\partial u}{\partial y}$, などが入り乱れた式
(\ao{偏微分方程式})になる
\end{itemize}
\end{frame}


%%%%%%%%%%%%%%%%% %%%%%%%%%%%%%%%%% 
\begin{frame}
\frametitle{物理で目にする例}
\begin{itemize}
\item 熱伝導: $u(x, y, z, t) =$ 場所$(x,y,z)$, 時刻$t$における温度

\[ c\frac{\partial \ao{u}}{\partial t}
= \lambda \left(\frac{\partial^2 \ao{u}}{\partial x^2}
+ \frac{\partial^2 \ao{u}}{\partial y^2}
+ \frac{\partial^2 \ao{u}}{\partial z^2}\right)
\]
\item 右辺の$(\cdots)$の部分は他の方程式でも何故か非常によく現れるので，
特別な記号($\Delta$または$\nabla^2$)が用意されている．

\[ \Delta u \equiv \frac{\partial^2 u}{\partial x^2}
+ \frac{\partial^2 u}{\partial y^2}
+ \frac{\partial^2 u}{\partial z^2}
\]

\item これを使えば，
\[ c\frac{\partial \ao{u}}{\partial t} = \Delta u \]
\end{itemize}


\end{frame}

%%%%%%%%%%%%%%%%% %%%%%%%%%%%%%%%%% 
\begin{frame}
\frametitle{物理で目にする例}
\begin{itemize}
\item マックスウェルの方程式(真空中)
  \begin{eqnarray*}
    \nabla \cdot \ao{B} & = & 0 \\
    \nabla \times \ao{E} & = & -\frac{\partial \ao{B}}{\partial t} \\
    \nabla \cdot \ao{E} & = & \frac{1}{\varepsilon}\rho \\
    \nabla \times \ao{B} & = & \mu j + \mu\varepsilon \ao{E} \\
  \end{eqnarray*}

注: $B$, $E$は\ao{ベクトル}場
  \begin{eqnarray*}
\nabla \cdot u & \equiv & \frac{\partial u_x}{\partial x}
+ \frac{\partial u_y}{\partial y}
+ \frac{\partial u_z}{\partial z} \\
\nabla \times u & \equiv & \ldots \mbox{教科書を見ま笑} \ldots \\
  \end{eqnarray*}
  
\item (一粒子の)シュレディンガーの方程式
\[ i\hbar \frac{\partial\ao{\psi}}{\partial t} 
= \left(\frac{-\hbar^2}{2m}\nabla^2 + V \right)\ao{\psi} \]
\end{itemize}

\end{frame}

%%%%%%%%%%%%%%%%% %%%%%%%%%%%%%%%%% 
\begin{frame}
\frametitle{物理で目にする例}
\begin{itemize}
\item 流体の速度場の方程式(ナビエストークス)
\[ 
\frac{\partial \ao{u}}{\partial t} + (\ao{u} \cdot \nabla)\ao{u}
= -\frac{1}{\rho} \nabla \ao{p} + \nu \nabla^2 \ao{u} + F
\]

\item 波動方程式(音とか光とか)
\[ \frac{1}{s^2} \frac{\partial^2 \ao{u}}{\partial t^2} = \Delta \ao{u} \] 

\end{itemize}
\end{frame}

%%%%%%%%%%%%%%%%% %%%%%%%%%%%%%%%%% 
\begin{frame}
\frametitle{コワイ??}

\begin{center}
\includegraphics[width=\textwidth]{out/pdf/svg/cat.pdf}
\end{center}

\end{frame}


%%%%%%%%%%%%%%%%% %%%%%%%%%%%%%%%%% 
\begin{frame}
\frametitle{しかし恐れる必要はない}
\begin{itemize}
\item 根本的な考えは同じ; (常)微分方程式が
\ao{たくさん並んでいるだけ}と思えば良い
\item 例えば1次元の場$u(x, t)$の方程式が，
\[ \frac{\partial u}{\partial t} = \ldots \]
と書けているならこれを用いて, 
{\footnotesize
\[
\left(
\begin{array}{c}
u(0,   t) \\
u(0.1,t) \\
u(0.2,t) \\
u(0.3,t) \\
u(0.4,t) \\
  \vdots  \\
\end{array}
\right)
\rightarrow
\left(
\begin{array}{c}
u(0,   t+\Delta t) \\
u(0.1,t+\Delta t) \\
u(0.2,t+\Delta t) \\
u(0.3,t+\Delta t) \\
u(0.4,t+\Delta t) \\
  \vdots  \\
\end{array}
\right)
\rightarrow
\left(
\begin{array}{c}
u(0,   t+2\Delta t) \\
u(0.1,t+2\Delta t) \\
u(0.2,t+2\Delta t) \\
u(0.3,t+2\Delta t) \\
u(0.4,t+2\Delta t) \\
  \vdots  \\
\end{array}
\right)
\rightarrow \cdots
\]}
という計算をしていくだけ(リストや配列が必須かつ有用)
\end{itemize}
\end{frame}


%%%%%%%%%%%%%%%%% %%%%%%%%%%%%%%%%% 
\begin{frame}
\frametitle{ともかく例を見る}
\begin{itemize}
\item [] 簡単のため1次元空間の例題($x$は単なる実数): 
\[\left\{
  \begin{array}{rcll}
{\displaystyle
\frac{\partial u}{\partial t}} & = & 3\, {\displaystyle \frac{\partial u}{\partial x}} & (0 < x < 1), \\
u(0, t) & = & 0 \\
u(1, t) & = & 0 \\
u(x, 0) & = & f(x)
  \end{array}\right.
\]

\item $u$は「場所($x$)」と「時刻($t$)」の関数

\item 1次元($x$が単なる実数)だから視覚的には，
弦が時間につれどう変化していくかを求める問題

\begin{center}
\def\svgwidth{\textwidth}
\input{out/tex/svg/wave.pdf_tex}
\end{center}
\end{itemize}
\end{frame}


%%%%%%%%%%%%%%%%% %%%%%%%%%%%%%%%%% 
\begin{frame}
\frametitle{離散化}
\begin{itemize}
\item 微分方程式:
\[ \frac{\partial u}{\partial t} = 3\, \frac{\partial u}{\partial x} \]

\item 離散化:
\[ \frac{u(x,t+\Delta t)-u(x,t)}{\Delta t} \approx 
3\, \frac{\partial u(x+\Delta x, t) - u(x, t)}{\Delta x} \]

\[ \therefore
u(x,t+\Delta t) \approx 
u(x,t) + (u(x+\Delta x, t) - u(x, t)) \frac{3\Delta t}{\Delta x} \]

\end{itemize}
\end{frame}

%%%%%%%%%%%%%%%%% %%%%%%%%%%%%%%%%% 
\begin{frame}
\frametitle{離散化}
\begin{itemize}
\item []
\[ \therefore
u(x,t+\Delta t) \approx 
u(x,t) + (u(x+\Delta x, t) - u(x, t)) \frac{3\Delta t}{\Delta x} \]

\begin{center}
\def\svgwidth{0.6\textwidth}
\input{out/tex/svg/wave_discretize.pdf_tex}
\end{center}

\item ある時点における自分と，右隣($+\Delta x$)の人の値がわかれば，
  自分の一瞬($+\Delta t$)後の値がわかる
\end{itemize}
\end{frame}

%%%%%%%%%%%%%%%%% %%%%%%%%%%%%%%%%% 
\begin{frame}
\frametitle{プログラム化}
\begin{itemize}
\item 用意する変数: 先の●たちを格納した一次元の配列 {\tt u}
\item {\tt u}の要素数を{\tt n}とする

\begin{center}
\def\svgwidth{0.8\textwidth}
\input{out/tex/svg/wave_discretize_u.pdf_tex}
\end{center}

\item これを踏まえ，更新式
\[ u(x,t+\Delta t) =
u(x,t) + (u(x+\Delta x, t) - u(x, t)) \frac{3\Delta t}{\Delta x} \]
に相当するPythonのプログラムは?
\end{itemize}
\end{frame}

%%%%%%%%%%%%%%%%% %%%%%%%%%%%%%%%%% 
\begin{frame}[fragile]
\frametitle{プログラム化}
\[ u(x,t+\Delta t) =
u(x,t) + (u(x+\Delta x, t) - u(x, t)) \frac{3\Delta t}{\Delta x} \]
\begin{itemize}
\item []
\begin{lstlisting}
for i in range(@\aka{\tt 1, n-1}@):
  u[i] = u[i] + (u[i+1] - u[i]) * (3*dt/dx)
\end{lstlisting}
\item ポイント: \aka{端の点({\tt u[0], u[n-1]})
はこの式では更新しない}．何故か?
\end{itemize}
\end{frame}

%%%%%%%%%%%%%%%%% %%%%%%%%%%%%%%%%% 
\begin{frame}[fragile]
\frametitle{端の点の扱い}
\begin{itemize}
\item 一般に端の点の扱いはよく考えて慎重に．
いい加減にやるとおかしな結果になる

\item まずプログラム上も，
右端の点({\tt u[n-1]})を更新しようとすると困ったことになる
\item []
\begin{lstlisting}[basicstyle=\ttfamily\footnotesize]
u[n-1] = u[n-1] + (@\aka{\tt u[n]}@ - u[n-1]) * (3*dt/dx)
\end{lstlisting}
\item {\tt u}は{\tt n}要素．\aka{{\tt u[n]}などという要素はない}

\item ただしそれが{\tt u[n-1]}
を更新しない\aka{理由}ではない

\item 微分方程式に立ち返る
\begin{eqnarray}
\frac{\partial u}{\partial t} & = & 3\, \frac{\partial u}{\partial x} \quad
(0 \aka{<} x \aka{<} 1) \label{eq:pde}
\end{eqnarray}
そもそも式(\ref{eq:pde})は，端点において成り立つとは言われていない
($0 < x < 1$でしか成り立たない式を$x = 0, 1$に適用するのは，
純粋に間違い)
\end{itemize}
\end{frame}

%%%%%%%%%%%%%%%%% %%%%%%%%%%%%%%%%% 
\begin{frame}[fragile]
\frametitle{端の点の扱い}
\begin{itemize}
\item 端の点は，別途与えられた条件で求める
\item 本問題の場合，$u(0,t) = u(1,t) = 0$だった
\item よって，
\begin{lstlisting}
for i in range(@\aka{\tt 1, n-1}@):
  u[i] = u[i] + (u[i+1] - u[i]) * (3*dt/dx)
u[0] = u[n-1] = 0
\end{lstlisting}
が与えられた条件に忠実な，正しい更新式
\item 本問においては, {\tt u[0], u[n-1]}は常に0なのだから，
実は毎回更新する必要もない
\begin{lstlisting}
for i in range(@\aka{\tt 1, n-1}@):
  u[i] += (u[i+1] - u[i]) * (3*dt/dx)
\end{lstlisting}
(注: {\tt A += B}は，{\tt A = A + B}と同じ意味)
\end{itemize}
\end{frame}


%%%%%%%%%%%%%%%%% %%%%%%%%%%%%%%%%% 
\begin{frame}[fragile]
\frametitle{numpyの配列機能を使ってスッキリと}
\begin{itemize}
\item for文を使った更新式:
\begin{lstlisting}
for i in range(@\aka{\tt 1, n-1}@):
  u[i] += (u[i+1] - u[i]) * (3*dt/dx)
\end{lstlisting}

\item 同じことはnumpyの配列に対する{\tt +, *}
などを使ってより簡潔に書ける
\begin{lstlisting}
u[1:n-1] += (u[2:n] - u[1:n-1]) * (3*dt/dx)
\end{lstlisting}

\item 簡潔なだけでなく，圧倒的に($>100$倍)速い
\item 10000要素の更新に必要な時間
  \begin{itemize}
  \item for文: $2 \times 10^{-2}$秒程度
  \item 配列で: $8 \times 10^{-5}$秒程度
  \end{itemize}
\end{itemize}

\end{frame}


\end{document}
